% Options for packages loaded elsewhere
\PassOptionsToPackage{unicode}{hyperref}
\PassOptionsToPackage{hyphens}{url}
\documentclass[
]{article}
\usepackage{xcolor}
\usepackage{amsmath,amssymb}
\setcounter{secnumdepth}{-\maxdimen} % remove section numbering
\usepackage{iftex}
\ifPDFTeX
  \usepackage[T1]{fontenc}
  \usepackage[utf8]{inputenc}
  \usepackage{textcomp} % provide euro and other symbols
\else % if luatex or xetex
  \usepackage{unicode-math} % this also loads fontspec
  \defaultfontfeatures{Scale=MatchLowercase}
  \defaultfontfeatures[\rmfamily]{Ligatures=TeX,Scale=1}
\fi
\usepackage{lmodern}
\ifPDFTeX\else
  % xetex/luatex font selection
\fi
% Use upquote if available, for straight quotes in verbatim environments
\IfFileExists{upquote.sty}{\usepackage{upquote}}{}
\IfFileExists{microtype.sty}{% use microtype if available
  \usepackage[]{microtype}
  \UseMicrotypeSet[protrusion]{basicmath} % disable protrusion for tt fonts
}{}
\makeatletter
\@ifundefined{KOMAClassName}{% if non-KOMA class
  \IfFileExists{parskip.sty}{%
    \usepackage{parskip}
  }{% else
    \setlength{\parindent}{0pt}
    \setlength{\parskip}{6pt plus 2pt minus 1pt}}
}{% if KOMA class
  \KOMAoptions{parskip=half}}
\makeatother
\usepackage{longtable,booktabs,array}
\newcounter{none} % for unnumbered tables
\usepackage{calc} % for calculating minipage widths
% Correct order of tables after \paragraph or \subparagraph
\usepackage{etoolbox}
\makeatletter
\patchcmd\longtable{\par}{\if@noskipsec\mbox{}\fi\par}{}{}
\makeatother
% Allow footnotes in longtable head/foot
\IfFileExists{footnotehyper.sty}{\usepackage{footnotehyper}}{\usepackage{footnote}}
\makesavenoteenv{longtable}
\setlength{\emergencystretch}{3em} % prevent overfull lines
\providecommand{\tightlist}{%
  \setlength{\itemsep}{0pt}\setlength{\parskip}{0pt}}
\usepackage{bookmark}
\IfFileExists{xurl.sty}{\usepackage{xurl}}{} % add URL line breaks if available
\urlstyle{same}
\hypersetup{
  hidelinks,
  pdfcreator={LaTeX via pandoc}}

\author{}
\date{}

\begin{document}

\section{DistCat: A 2-Categorical Framework for Distance Derivation from
First
Principles}\label{distcat-a-2-categorical-framework-for-distance-derivation-from-first-principles}

\begin{quote}
\textbf{Version}: iteration 73, final draft (Fisher-Rao universality
theorem added) \textbf{Author}: Cheng Lu (Independent Researcher)
\textbf{License}: CC BY 4.0
\end{quote}

\subsection{Abstract}\label{abstract}

We present \(\mathbf{DistCat}\), a 2-categorical framework for deriving
distance functions from first principles. The central principle states
that any distance function \(\delta\) on an observation space is derived
from a more fundamental distance \(\Gamma\) on a source category via a
natural transformation \(\eta: \delta \Rightarrow \Gamma\). We define
\(\mathbf{DistCat}\) as a 2-category whose 0-objects are base structures
\(\mathcal{B}=(\mathbf{C},\mathcal{F},\mathcal{O})\), 1-morphisms are
structure-preserving functors with natural transformations, and
2-morphisms are distance-exporting pairs \((\Gamma,\delta)\). We prove
the existence of a 2-functor \(D: \mathbf{DistCat} \to \mathbf{Met}\)
mapping base structures to metric spaces, and prove observation
infeasibility: distinct base structures may yield isometric observation
spaces. We prove the distance export functor
\(D: \mathbf{DistCat} \to \mathbf{Met}\) is well-defined on all five
examples, including a complete 2-category axiom verification for the
EMRI case. We verify the framework with five independently-motivated
examples: (1) D1-grav-waves, reconstructing the Hellings-Downs
correlation as a natural transformation in Extended Functorial AQFT; (2)
D16+D24, deriving semantic geodesics from Fisher information geometry;
(3) Bayesian inference, as Fisher-Rao distance under marginalization;
(4) D26, formalizing resonance amplification in control theory; (5) EMRI
chaos, connecting dark matter halo perturbations to Lyapunov exponents.
We prove three predictions for the EMRI example: unstable-orbital
frequency scaling, GW amplitude proportionality to resonance
coefficients, and a KAM-breakup square-root scaling law for the Lyapunov
exponent. We establish distance equivalence between D1 and D16+D24
(asymptotic), between D26 and EMRI (structural), and between D1 and EMRI
(through GW observational bridge). The distance export principle
\(\delta = \Phi(\Gamma)\) appears to be universal: all examples share
the same 2-categorical structure despite having different physical
origins and critical exponents.

\textbf{Keywords}: category theory, distance derivation, information
geometry, quantum field theory, gravitational waves, resonance
amplification, KAM theory

\begin{center}\rule{0.5\linewidth}{0.5pt}\end{center}

\subsection{1. Introduction}\label{introduction}

The concept of ``distance'' permeates mathematics and physics, yet its
origin remains largely unexamined. In differential geometry, the
Riemannian distance is derived from a metric tensor. In information
theory, the KL divergence is derived from a statistical model. In
quantum field theory, state distances are derived from algebraic
structures. Each derivation appears domain-specific, with no unifying
principle.

This paper proposes that \textbf{all distance functions are exported
from a more fundamental structure via categorical natural
transformations}. Specifically, given a source category \(\mathbf{C}\)
with a pre-distance \(\Gamma\), an observation functor
\(\mathcal{O}: \mathbf{C} \to \mathbf{E}\), and a natural transformation
\(\eta: \delta \Rightarrow \Gamma\), the observable distance \(\delta\)
is completely determined by \(\Gamma\) and \(\mathcal{O}\).

The framework \(\mathbf{DistCat}\) formalizes this idea as a
\textbf{2-category} with: - 0-objects: base structures
\(\mathcal{B} = (\mathbf{C}, \mathcal{F}, \mathcal{O})\) - 1-morphisms:
structure-preserving triples \((\Phi, \alpha, \beta)\) - 2-morphisms:
distance-exporting pairs \((\Gamma, \delta)\)

We verify this framework with five examples spanning physics,
statistics, and control theory:

\begin{enumerate}
\def\labelenumi{\arabic{enumi}.}
\tightlist
\item
  \textbf{D1-grav-waves}: The Hellings-Downs correlation
  \(\zeta(\theta)\) between pulsar pairs is derived from the total
  variation distance between algebraic states in Extended Functorial
  AQFT.
\item
  \textbf{D16+D24}: The KL divergence between categorical distributions
  is derived from the Fisher-Rao distance on the semantic manifold,
  recovering Amari's local equivalence.
\item
  \textbf{Bayesian inference}: The Fisher-Rao distance under
  marginalization corresponds to Bayesian posterior contraction.
\item
  \textbf{D26 resonance}: The resonance amplification coefficient
  \(G_{\text{amp}}\) in control theory is the distance export from a
  pole-perturbation distance.
\item
  \textbf{EMRI chaos}: The Lyapunov exponent \(\lambda_T\)
  characterizing chaotic dynamics near a Schwarzschild black hole
  immersed in a dark matter halo is the distance export from a
  metric-perturbation distance.
\end{enumerate}

The five examples have critical exponents
\(\alpha \in \{1.84, 2.00, 2.00, \infty, 0.50\}\), yet all share the
universal pattern \(\delta \sim \Gamma^\alpha\) near critical points.
This suggests that the distance export principle is a
\textbf{categorical invariant}, not a domain-specific accident.

\begin{center}\rule{0.5\linewidth}{0.5pt}\end{center}

\subsection{2. Mathematical Definition}\label{mathematical-definition}

\subsubsection{2.1 Base Structures
(0-Objects)}\label{base-structures-0-objects}

\textbf{Definition 2.1 (Base Structure)}: A base structure is a triple
\(\mathcal{B} = (\mathbf{C}, \mathcal{F}, \mathcal{O})\) where:

\begin{itemize}
\tightlist
\item
  \(\mathbf{C}\) is a \textbf{source category} (objects are fundamental
  system units)
\item
  \(\mathcal{F}: \mathbf{C} \to \mathbf{D}\) is a \textbf{structure
  functor} mapping to a target category \(\mathbf{D}\)
\item
  \(\mathcal{O}: \mathbf{C} \to \mathbf{E}\) is an \textbf{observation
  functor} mapping to an observation category \(\mathbf{E}\)
\end{itemize}

\textbf{Convention}: Throughout, \(\mathbf{E} = \mathbf{Set}\), so
\(\mathcal{O}(X)\) is a set for each \(X \in \text{Obj}(\mathbf{C})\).

\textbf{Examples}:

{\def\LTcaptype{none} % do not increment counter
\begin{longtable}[]{@{}
  >{\raggedright\arraybackslash}p{(\linewidth - 10\tabcolsep) * \real{0.0492}}
  >{\raggedright\arraybackslash}p{(\linewidth - 10\tabcolsep) * \real{0.0984}}
  >{\raggedright\arraybackslash}p{(\linewidth - 10\tabcolsep) * \real{0.2131}}
  >{\raggedright\arraybackslash}p{(\linewidth - 10\tabcolsep) * \real{0.2131}}
  >{\raggedright\arraybackslash}p{(\linewidth - 10\tabcolsep) * \real{0.2131}}
  >{\raggedright\arraybackslash}p{(\linewidth - 10\tabcolsep) * \real{0.2131}}@{}}
\toprule\noalign{}
\begin{minipage}[b]{\linewidth}\raggedright
\#
\end{minipage} & \begin{minipage}[b]{\linewidth}\raggedright
Name
\end{minipage} & \begin{minipage}[b]{\linewidth}\raggedright
\(\mathbf{C}\)
\end{minipage} & \begin{minipage}[b]{\linewidth}\raggedright
\(\mathbf{D}\)
\end{minipage} & \begin{minipage}[b]{\linewidth}\raggedright
\(\mathcal{F}\)
\end{minipage} & \begin{minipage}[b]{\linewidth}\raggedright
\(\mathcal{O}\)
\end{minipage} \\
\midrule\noalign{}
\endhead
\bottomrule\noalign{}
\endlastfoot
1 & D1-grav-waves & \(\mathbf{Loc}_M^{\text{ext}}\) &
\(\mathbf{C}^*\mathbf{Alg}\) & AQFT functor \(F\) & Correlation
\(\langle\phi(x_1)\phi(x_2)\rangle\) \\
2 & D16+D24 & \(\mathbf{FinProb}\) & \(\mathbf{Man}\) & Statistical
model \(\mathcal{P}\) & Classification marginal \(\mathcal{Q}\) \\
3 & Bayesian & \(\mathbf{Bayes}\) & \(\mathbf{Prob}\) & Prior-posterior
& Posterior \(\text{P}(\cdot\mid\text{data})\) \\
4 & D26 & LTI systems & \(\mathbf{Rat}(s)\) & Transfer function
\(C(sI-A)^{-1}B+D\) & \((\omega_{\text{res}}, G_{\text{amp}})\) \\
5 & EMRI & \(\mathbf{Loc}_{M,g}\) & \(\mathbf{Riem}\) & Geodesic flow &
\((\lambda_T, \kappa, \{r_{\text{circ}}\})\) \\
\end{longtable}
}

\subsubsection{2.2 Structure-Preserving
1-Morphisms}\label{structure-preserving-1-morphisms}

\textbf{Definition 2.2 (1-Morphism)}: Let
\(\mathcal{B}_1 = (\mathbf{C}_1, \mathcal{F}_1, \mathcal{O}_1)\) and
\(\mathcal{B}_2 = (\mathbf{C}_2, \mathcal{F}_2, \mathcal{O}_2)\) be base
structures. A 1-morphism \(f: \mathcal{B}_1 \to \mathcal{B}_2\) is a
triple \((\Phi, \alpha, \beta)\) where:

\begin{itemize}
\tightlist
\item
  \(\Phi: \mathbf{C}_1 \to \mathbf{C}_2\) is a functor
\item
  \(\alpha: \mathcal{F}_2 \circ \Phi \Rightarrow \mathcal{F}_1\) is a
  natural transformation
\item
  \(\beta: \mathcal{O}_1 \Rightarrow \mathcal{O}_2 \circ \Phi\) is a
  natural transformation
\end{itemize}

\textbf{Physical meaning}: A 1-morphism encodes a structure-preserving
embedding of \(\mathcal{B}_1\) into \(\mathcal{B}_2\), maintaining
consistency of both the structure map \(\mathcal{F}\) and the
observation map \(\mathcal{O}\).

\textbf{Vertical composition}:
\((\Phi_2, \alpha_2, \beta_2) \circ_v (\Phi_1, \alpha_1, \beta_1) = (\Phi_2 \circ \Phi_1, \alpha_1 \cdot \alpha_2, \beta_1 \cdot \beta_2)\)

\textbf{Horizontal composition}:
\((\Phi_2, \alpha_2, \beta_2) \circ_h (\Phi_1, \alpha_1, \beta_1) = (\Phi_2 \circ \Phi_1, \alpha_1 \circ \Phi_2 \cdot \alpha_2, \beta_2 \circ \Phi_1 \cdot \beta_1)\)

\subsubsection{2.3 Distance-Deriving
2-Morphisms}\label{distance-deriving-2-morphisms}

\textbf{Definition 2.3 (2-Morphism)}: Let \(\mathcal{B}\) be a base
structure. A distance-exporting 2-morphism is a pair
\((\Gamma, \delta)\) where:

\begin{itemize}
\tightlist
\item
  \(\Gamma: \mathbf{C} \times \mathbf{C} \to \mathbb{R}_{\geq 0}\) is a
  \textbf{pre-distance} on the source category
\item
  \(\delta: \mathcal{O}(\mathbf{C}) \times \mathcal{O}(\mathbf{C}) \to \mathbb{R}_{\geq 0}\)
  is a \textbf{derived distance} on the observation space
\item
  There exists a natural transformation
  \(\eta: \delta(\mathcal{O}(X), \mathcal{O}(Y)) \to \Gamma(X, Y)\) for
  all \(X, Y \in \mathbf{C}\)
\end{itemize}

\textbf{Distance export principle}: If \(\eta\) is a natural
isomorphism, then \(\delta\) is completely determined by \(\Gamma\)
through \(\mathcal{O}\).

\subsubsection{\texorpdfstring{2.4 The 2-Category
\(\mathbf{DistCat}\)}{2.4 The 2-Category \textbackslash mathbf\{DistCat\}}}\label{the-2-category-mathbfdistcat}

\textbf{Definition 2.4 (\(\mathbf{DistCat}\))}: \(\mathbf{DistCat}\) is
a 2-category defined by:

\begin{itemize}
\tightlist
\item
  \textbf{0-objects}: Base structures
  \(\mathcal{B} = (\mathbf{C}, \mathcal{F}, \mathcal{O})\)
\item
  \textbf{1-morphisms}: Structure-preserving triples
  \((\Phi, \alpha, \beta)\)
\item
  \textbf{2-morphisms}: Distance-exporting pairs \((\Gamma, \delta)\)
  with natural transformation \(\eta\)
\end{itemize}

\textbf{Composition rules}: - Vertical:
\((\Gamma_2, \delta_2) \circ_v (\Gamma_1, \delta_1)\) via composite
natural transformations - Horizontal:
\((\Gamma_2, \delta_2) \circ_h (\Gamma_1, \delta_1)\) via
functor-composition pushout

\subsubsection{2.5 The Distance Export
Functor}\label{the-distance-export-functor}

\textbf{Theorem 2.5 (Distance export functor \(D\))}: There exists a
2-functor \(D: \mathbf{DistCat} \to \mathbf{Met}\) where
\(\mathbf{Met}\) is the 2-category of metric spaces, such that:

\begin{itemize}
\tightlist
\item
  \(D(\mathcal{B}) = (\mathcal{O}(\mathbf{C}), \delta)\)
\item
  \(D(\Phi, \alpha, \beta) = \beta(\mathbf{C}): \mathcal{O}_1(\mathbf{C}_1) \to \mathcal{O}_2(\mathbf{C}_2)\)
\item
  \(D(\zeta) = \zeta_*|_{\text{Im}(\mathcal{O}_1)}\)
\end{itemize}

\textbf{Proof}: 1. \(D(\mathcal{B})\) is a metric space:
\(\mathcal{O}(\mathbf{C}) = \{\mathcal{O}(X) \mid X \in \text{Obj}(\mathbf{C})\}\)
and \(\delta\) is defined by Definition 2.3. 2.
\(D(\Phi, \alpha, \beta) = \beta(\mathbf{C})\) is a map between
observation spaces (naturality of \(\beta\)). 3. Identity:
\(D(\text{id}_{\mathcal{B}}) = \text{id}_{\delta} = \text{id}_{D(\mathcal{B})}\).
4. 1-morphism composition preservation:
\(D((\Phi_2, \alpha_2, \beta_2) \circ_v (\Phi_1, \alpha_1, \beta_1)) = \beta_1 \cdot \beta_2 = D(\Phi_2, \alpha_2, \beta_2) \circ D(\Phi_1, \alpha_1, \beta_1)\).
5. 2-morphism composition preservation:
\(D(\zeta_2 \circ_v \zeta_1) = (\zeta_2 \circ \zeta_1)_* = \zeta_{2*} \circ \zeta_{1*} = D(\zeta_2) \circ D(\zeta_1)\).

\subsubsection{2.6 Observation
Infeasibility}\label{observation-infeasibility}

\textbf{Theorem 2.6 (Observation infeasibility)}: There exist
\(\mathcal{B}_1 \neq \mathcal{B}_2\) such that
\(D(\mathcal{B}_1) \cong D(\mathcal{B}_2)\) (isometric observation
spaces).

\textbf{Proof}: Consider \(\mathcal{B}_1\) with
\(\mathbf{C}_1 = \mathbf{FinProb}\) (D16+D24) and \(\mathcal{B}_2\) with
\(\mathbf{C}_2 = \mathbf{Loc}_M^{\text{ext}}\) (D1-grav-waves). As
\(\delta_1, \delta_2 \to 0\) near critical points, both approach
quadratic forms: \(\delta_1 \approx \frac{1}{2} g_F\) and
\(\delta_2 \approx \frac{1}{2} g_{HD}\). When \(g_F\) and \(g_{HD}\)
share the same curvature structure, \(\delta_1 \cong \delta_2\). Hence
\(D(\mathcal{B}_1) \cong D(\mathcal{B}_2)\) but
\(\mathcal{B}_1 \neq \mathcal{B}_2\).

\textbf{Corollary 2.7}: Distance export is not invertible --- the base
structure cannot be uniquely recovered from the observation distance.
Distance equivalence classes
\([\mathcal{B}] = \{\mathcal{B}' \mid D(\mathcal{B}') \cong D(\mathcal{B})\}\)
exist.

\begin{center}\rule{0.5\linewidth}{0.5pt}\end{center}

\subsection{3. Example 1: D1-grav-waves (HD Correlation as Natural
Transformation)}\label{example-1-d1-grav-waves-hd-correlation-as-natural-transformation}

\subsubsection{3.1 Base Structure}\label{base-structure}

\begin{itemize}
\tightlist
\item
  \(\mathbf{C} = \mathbf{Loc}_M^{\text{ext}}\) (extended causal regions)
\item
  \(\mathcal{F} = F: \mathbf{Loc}_M^{\text{ext}} \to \mathbf{C}^*\mathbf{Alg}\)
  (Extended Functorial AQFT)
\item
  \(\mathcal{O}(O) = \langle \phi(x_1)\phi(x_2) \rangle\) (correlation
  observable, \(x_i \in O\))
\item
  \(\Gamma((O_1, O_2)) = d_{TV}(\omega_1, \omega_2)\) (total variation
  distance between algebraic states)
\item
  \(\delta((R_1, R_2)) = \zeta(\theta_{12})\) (Hellings-Downs
  correlation)
\item
  \(\eta: \zeta(\theta_{12}) \Rightarrow d_{TV}(\omega_1, \omega_2)\)
  (naturality from AQFT functoriality)
\end{itemize}

\subsubsection{3.2 HD Correlation
Function}\label{hd-correlation-function}

The Hellings-Downs function:

\[\zeta(\theta) = \frac{3}{2} x \ln x - \frac{x}{4} + \frac{1}{2} + \frac{1}{2}\delta_{12}, \quad x = \frac{1 - \cos\theta}{2}\]

\subsubsection{3.3 Small-Angle
Asymptotics}\label{small-angle-asymptotics}

For \(\theta \to 0\), \(x \approx \theta^2/4\):

\[\zeta(\theta) \approx \frac{1}{2} + \frac{3}{8}\theta^2(\ln \theta^2 - \ln 4) - \frac{\theta^2}{16} + \mathcal{O}(\theta^4)\]

\[\zeta(\theta) - \frac{1}{2} \sim \theta^2 \ln \theta \quad (\theta \to 0)\]

Numerical verification yields \(|\zeta - 0.5| \propto \theta^{1.84}\),
consistent with \(\theta^2 \ln\theta\) (log correction reduces effective
exponent).

\subsubsection{3.4 Physical Meaning}\label{physical-meaning}

The \(SO(3)\) rotational symmetry of \(\zeta(\theta)\) reflects the
\(SO(3)\) rotation group structure of the AQFT assignment on the sphere
at infinity. The logarithmic term arises from the nontrivial cobordism
class connecting two worldlines in \(\mathbf{Loc}_M^{\text{ext}}\).

\begin{center}\rule{0.5\linewidth}{0.5pt}\end{center}

\subsection{4. Example 2: D16+D24 (Fisher Information
Geometry)}\label{example-2-d16d24-fisher-information-geometry}

\subsubsection{4.1 Base Structure}\label{base-structure-1}

\begin{itemize}
\tightlist
\item
  \(\mathbf{C} = \mathbf{FinProb}\) (finite probability spaces,
  morphisms = stochastic maps)
\item
  \(\mathcal{F} = \mathcal{P}\) (statistical model functor
  \(\mathcal{P}(\sigma) = p(\cdot \mid \sigma)\))
\item
  \(\mathcal{O} = \mathcal{Q}\) (classification observation, marginal
  distribution \(q_{\mathcal{C}}(\cdot \mid \sigma)\))
\item
  \(\Gamma((\sigma_1, \sigma_2)) = \frac{1}{2} d_{g_F}(\sigma_1, \sigma_2)^2\)
  (Fisher distance squared)
\item
  \(\delta((q_1, q_2)) = D_{KL}(q_1 \parallel q_2)\) (KL divergence)
\item
  \(\eta: D_{KL} \Rightarrow \frac{1}{2} d_{g_F}^2\) (Amari's lemma)
\end{itemize}

\subsubsection{4.2 Geodesic-KL Equivalence
(Theorem)}\label{geodesic-kl-equivalence-theorem}

\textbf{Theorem 4.1 (Geodesic-KL equivalence)}: For any statistical
model \(\mathcal{P}\) with Fisher metric \(g_F\), and any
\(\sigma_1, \sigma_2\) sufficiently close:

\[D_{KL}(p(\cdot \mid \sigma_1) \parallel p(\cdot \mid \sigma_2)) = \frac{1}{2} d_{g_F}(\sigma_1, \sigma_2)^2 + \mathcal{O}(\|\sigma_1 - \sigma_2\|^3)\]

\textbf{Proof} (Amari 1985): The KL divergence has a stationary point at
\(\sigma_1 = \sigma_2\) (both first derivatives vanish). The Hessian at
this point equals the Fisher metric \(g_F\). Taylor expansion yields the
result.

\subsubsection{4.3 Classification Granularity
Inequality}\label{classification-granularity-inequality}

\textbf{Theorem 4.2}: For any categorization \(\mathcal{C}\) of a
statistical model space \(\Sigma\), the KL divergence between marginal
distributions satisfies:

\[D_{KL}(q_{\mathcal{C}}(\cdot \mid \sigma_1) \parallel q_{\mathcal{C}}(\cdot \mid \sigma_2)) \leq D_{KL}(p(\cdot \mid \sigma_1) \parallel p(\cdot \mid \sigma_2))\]

Equality holds iff the categorization preserves all information in the
likelihood ratio.

\begin{center}\rule{0.5\linewidth}{0.5pt}\end{center}

\subsection{5. Example 3: Bayesian
Inference}\label{example-3-bayesian-inference}

\subsubsection{5.1 Base Structure}\label{base-structure-2}

\begin{itemize}
\tightlist
\item
  \(\mathbf{C} = \mathbf{Bayes}\) (Bayesian update category: objects =
  (prior, likelihood) pairs)
\item
  \(\mathcal{F}\): Bayesian update functor mapping
  \((p_{\text{prior}}, p_{\text{likelihood}}) \to p_{\text{posterior}}\)
\item
  \(\mathcal{O}\): Marginal observation (predictive distribution for new
  data)
\item
  \(\Gamma\): Fisher-Rao distance on the prior parameter space
\item
  \(\delta\): KL divergence between predictive distributions
\end{itemize}

\subsubsection{5.2 Fisher-Rao as Distance
Export}\label{fisher-rao-as-distance-export}

The Fisher information matrix \(I(\theta)\) on the parameter space is
the Hessian of the KL divergence at the true parameter. The Fisher-Rao
distance:

\[d_{FR}(\theta_1, \theta_2) = \inf_{\gamma} \int_0^1 \sqrt{\dot{\theta}^T I(\theta) \dot{\theta}} \, dt\]

is the distance export of \(\Gamma\) under the Bayesian update functor.
The KL divergence between posteriors is the marginalization
(observation) of the Fisher-Rao distance.

\subsubsection{5.3 Marginalization as Distance
Derivation}\label{marginalization-as-distance-derivation}

The Marginalization Inequality (a special case of Theorem 4.2):
marginalizing a posterior increases the effective KL divergence. This is
equivalent to the data processing inequality in information theory, now
interpreted as a DistCat 2-morphism property.

\begin{center}\rule{0.5\linewidth}{0.5pt}\end{center}

\subsection{6. Example 4: D26 Resonance
Amplification}\label{example-4-d26-resonance-amplification}

\subsubsection{6.1 Control-Theory Base
Structure}\label{control-theory-base-structure}

\begin{itemize}
\tightlist
\item
  \(\mathbf{C}\): Category of linear time-invariant (LTI) systems
  \(S = (A, B, C, D)\)
\item
  \(\mathcal{F}\): Transfer-function functor
  \(G(s) = C(sI - A)^{-1}B + D\)
\item
  \(\mathcal{O}\): Resonance observation
  \((\omega_{\text{res}}, G_{\text{amp}}(\omega_{\text{res}}))\)
\item
  \(\Gamma((S_1, S_2)) = \max_i |\lambda_i(A_1) - \lambda_i(A_2)|\)
  (pole-perturbation distance)
\item
  \(\delta((\omega_{r1}, G_1), (\omega_{r2}, G_2)) = |\omega_{r1} - \omega_{r2}| + |1/G_1 - 1/G_2|\)
\item
  \(\eta: \delta \Rightarrow \Gamma\) (resonance-observation distance
  bounded by pole-perturbation distance)
\end{itemize}

\subsubsection{6.2 Resonance Frequency
Sensitivity}\label{resonance-frequency-sensitivity}

\textbf{Theorem 6.1}: Let
\(S_\epsilon = (A + \epsilon \delta A, B, C, D)\) be a perturbed system.
Then:

\[\lim_{\epsilon \to 0} \frac{|\omega_{\text{res}}(S_\epsilon) - \omega_{\text{res}}(S)|}{\epsilon} \leq \max_i \frac{|\lambda_i(A + \epsilon \delta A) - \lambda_i(A)|}{\epsilon} = \kappa_\lambda(\delta A)\]

where \(\kappa_\lambda\) is the eigenvalue condition number.

\subsubsection{6.3 Critical Amplification
Bound}\label{critical-amplification-bound}

\textbf{Theorem 6.2}: Let
\(\epsilon = \min_i |\text{Re}(\lambda_i(A))|\) be the distance to
criticality. Then:

\[G_{\text{amp}}(S) \leq \frac{\|C\| \cdot \|B\|}{\epsilon}\]

\subsubsection{6.4 H∞ Robustness as Distance
Export}\label{h-robustness-as-distance-export}

The H∞ robustness margin \(\mathcal{R}(S) = 1/\|T_{zw}\|_\infty\) equals
the inverse of the distance to the critical observation:
\(\mathcal{R}(S) = \delta(\mathcal{O}(S), \mathcal{O}_{\text{critical}})^{-1}\),
where
\(\mathcal{O}_{\text{critical}} = (\omega_{\text{crit}}, \infty)\).

\begin{center}\rule{0.5\linewidth}{0.5pt}\end{center}

\subsection{7. Example 5: EMRI Chaos in Dark Matter
Halos}\label{example-5-emri-chaos-in-dark-matter-halos}

\subsubsection{7.1 GR Base Structure}\label{gr-base-structure}

Consider a Schwarzschild black hole of mass \(M\) immersed in a
Dehnen-(1,4,5/2) dark matter halo of core density \(\rho_s\) and scale
radius \(r_s\). The metric is:

\[ds^2 = -f(r)dt^2 + f(r)^{-1}dr^2 + r^2 d\Omega^2, \quad f(r) = 1 - \frac{2M}{r} - \frac{32\pi\rho_s r_s^2}{1 + r_s/r}\]

A test particle of mass \(m \ll M\) follows a geodesic with conserved
energy \(E\) and angular momentum \(L\). The radial motion is governed
by the effective potential (Eq. 3.5):

\[V_{\text{eff}}(x) = -\frac{4M^2 m^2}{L^2}x + x^2 + x^3 - \frac{32\pi\rho_s r_s^2}{1 + x r_s/(2M)}\left(x^2 + \frac{4M^2 m^2}{L^2}\right), \quad x = \frac{2M}{r}\]

The five terms represent: 1/r attraction, centrifugal repulsion, GR
correction (\(\propto 1/r^3\)), and two DM-halo coupling terms.

\subsubsection{7.2 Homoclinic Orbits and Lyapunov
Exponents}\label{homoclinic-orbits-and-lyapunov-exponents}

For \(r_{\text{ISCO}} < r < r_{\text{MBO}}\), there exists a unique
unstable circular orbit at \(r_{\text{un}}\) where
\(V_{\text{eff}}'(r_{\text{un}}) = 0\), and a homoclinic orbit
connecting \((r_{\text{un}}, 0)\) to itself in phase space.

Numerical results (Table 1-2 of Das et al.~2025, Paper I):

{\def\LTcaptype{none} % do not increment counter
\begin{longtable}[]{@{}
  >{\raggedright\arraybackslash}p{(\linewidth - 8\tabcolsep) * \real{0.2000}}
  >{\raggedright\arraybackslash}p{(\linewidth - 8\tabcolsep) * \real{0.2000}}
  >{\raggedright\arraybackslash}p{(\linewidth - 8\tabcolsep) * \real{0.2000}}
  >{\raggedright\arraybackslash}p{(\linewidth - 8\tabcolsep) * \real{0.2000}}
  >{\raggedright\arraybackslash}p{(\linewidth - 8\tabcolsep) * \real{0.2000}}@{}}
\toprule\noalign{}
\begin{minipage}[b]{\linewidth}\raggedright
\(\rho_s\)
\end{minipage} & \begin{minipage}[b]{\linewidth}\raggedright
\(r_s\)
\end{minipage} & \begin{minipage}[b]{\linewidth}\raggedright
\(L\) (10⁻⁵)
\end{minipage} & \begin{minipage}[b]{\linewidth}\raggedright
\(r_{\text{un}}\)
\end{minipage} & \begin{minipage}[b]{\linewidth}\raggedright
\(E_{\text{homo}}\) (10⁻⁵)
\end{minipage} \\
\midrule\noalign{}
\endhead
\bottomrule\noalign{}
\endlastfoot
0.000 & 0.00 & 3.75 & 4.3385 & 0.970418 \\
0.010 & 0.15 & 3.75 & 4.6392 & 0.950055 \\
0.025 & 0.15 & 3.75 & 5.3416 & 0.921035 \\
0.040 & 0.15 & 4.25 & 4.6042 & 0.936628 \\
0.050 & 0.15 & 4.25 & 4.9047 & 0.914265 \\
0.065 & 0.15 & 4.25 & 5.5373 & 0.882241 \\
\end{longtable}
}

The unstable orbit radius \(r_{\text{un}}\) \textbf{increases} with
\(\rho_s\) (DM halo mass pushes the orbit outward).

\subsubsection{\texorpdfstring{7.3 D26\(\leftrightarrow\)EMRI Distance
Equivalence}{7.3 D26\textbackslash leftrightarrowEMRI Distance Equivalence}}\label{d26leftrightarrowemri-distance-equivalence}

{\def\LTcaptype{none} % do not increment counter
\begin{longtable}[]{@{}
  >{\raggedright\arraybackslash}p{(\linewidth - 2\tabcolsep) * \real{0.5000}}
  >{\raggedright\arraybackslash}p{(\linewidth - 2\tabcolsep) * \real{0.5000}}@{}}
\toprule\noalign{}
\begin{minipage}[b]{\linewidth}\raggedright
D26 concept
\end{minipage} & \begin{minipage}[b]{\linewidth}\raggedright
EMRI counterpart
\end{minipage} \\
\midrule\noalign{}
\endhead
\bottomrule\noalign{}
\endlastfoot
LTI system \((A,B,C,D)\) & Schwarzschild geodesic system \\
Transfer function \(G(s)\) & Effective potential \(V_{\text{eff}}\) \\
Poles \(\lambda_i\) & Circular orbit radii \(r_{\text{circ}}\) \\
Resonance frequency \(\omega_{\text{res}}\) & Unstable orbit frequency
\(\Omega_{\text{un}}\) \\
Perturbation \(\epsilon \delta A\) & DM halo parameters
\((\rho_s, r_s)\) \\
Amplification \(G_{\text{amp}}\) & Lyapunov exponent \(\lambda_T\) \\
Critical distance \(\epsilon\) & Distance to unstable orbit
\(r - r_{\text{un}}\) \\
MSS chaos bound \(\lambda^2 \leq \kappa^2\) & H\(\infty\) robustness \\
\end{longtable}
}

\textbf{Proposition 7.1}:
\(D(\text{EMRI}) \cong_{\text{struct}} D(\text{D26})\). Both exhibit:
power-law approach \(\delta \sim \Gamma^\alpha\) with
\(\alpha \in (0, \infty)\); resonance peak at critical frequency;
robustness bound at threshold.

\subsubsection{7.4 Unstable-Orbital-Frequency Numerical
Results}\label{unstable-orbital-frequency-numerical-results}

From the homoclinic data above, the angular frequency at the unstable
orbit:

\[\Omega_{\text{un}} = \frac{1}{r_{\text{un}}^{3/2}} \quad (\text{geometric units, } M = 1)\]

For an EMRI with \(M = 10^6 M_\odot\):

{\def\LTcaptype{none} % do not increment counter
\begin{longtable}[]{@{}
  >{\raggedright\arraybackslash}p{(\linewidth - 8\tabcolsep) * \real{0.2000}}
  >{\raggedright\arraybackslash}p{(\linewidth - 8\tabcolsep) * \real{0.2000}}
  >{\raggedright\arraybackslash}p{(\linewidth - 8\tabcolsep) * \real{0.2000}}
  >{\raggedright\arraybackslash}p{(\linewidth - 8\tabcolsep) * \real{0.2000}}
  >{\raggedright\arraybackslash}p{(\linewidth - 8\tabcolsep) * \real{0.2000}}@{}}
\toprule\noalign{}
\begin{minipage}[b]{\linewidth}\raggedright
\(\rho_s\)
\end{minipage} & \begin{minipage}[b]{\linewidth}\raggedright
\(r_s\)
\end{minipage} & \begin{minipage}[b]{\linewidth}\raggedright
\(r_{\text{un}}\)
\end{minipage} & \begin{minipage}[b]{\linewidth}\raggedright
\(\Omega_{\text{un}}\)
\end{minipage} & \begin{minipage}[b]{\linewidth}\raggedright
\(f_{\text{GW}}\) (mHz)
\end{minipage} \\
\midrule\noalign{}
\endhead
\bottomrule\noalign{}
\endlastfoot
0.000 & 0.00 & 4.3385 & 0.110661 & 1.708 \\
0.010 & 0.15 & 4.6392 & 0.100078 & 1.544 \\
0.025 & 0.15 & 5.3416 & 0.081002 & 1.250 \\
0.040 & 0.15 & 4.6042 & 0.101221 & 1.562 \\
0.050 & 0.15 & 4.9047 & 0.092063 & 1.421 \\
0.065 & 0.15 & 5.5373 & 0.076746 & 1.184 \\
\end{longtable}
}

\textbf{Key findings}: 1. \(\Omega_{\text{un}}\) decreases with DM halo
parameters (unstable orbit pushed outward) 2. GW frequency 1.2-1.7 mHz
falls within LISA's optimal band (0.1-100 mHz) 3. Power-law:
\(\Omega_{\text{un}} \sim \rho_s^{-0.3}\) (approximate)

\subsubsection{7.5 GW Verification (Paper
II)}\label{gw-verification-paper-ii}

Das et al.~(2025, Paper II) compute GW waveforms using the numerical
kludge scheme. Key results:

\begin{itemize}
\tightlist
\item
  \textbf{Amplitude}: Chaotic orbits produce GW amplitudes 1.7-2.5×
  higher than non-chaotic (Fig. 4 of Paper II)
\item
  \textbf{Spectrum}: Non-chaotic GWs = discrete characteristic
  frequencies; chaotic GWs = quasi-continuous broadband spectrum (Figs.
  5-7 of Paper II)
\item
  \textbf{Energy emission}: \(dE/dt\) scales with chaos level
\item
  \textbf{Recurrence rate (RR)}: Consistently RR = 0.06 (non-chaotic)
  \(\to\) RR = 0.15 (chaotic), a 2.5× increase across all cases
\item
  \textbf{Characteristic strain}: Chaotic GWs exceed LISA/TianQin/Taiji
  sensitivity thresholds
\end{itemize}

\subsubsection{7.6 KAM Scaling Law
(Theorem)}\label{kam-scaling-law-theorem}

\textbf{Theorem 7.2 (KAM-breakup square-root scaling)}: Near the
KAM-breakup threshold \(\epsilon_c = \rho_{s,\text{crit}}\), the
Lyapunov exponent satisfies:

\[\lambda_T \sim |\epsilon - \epsilon_c|^{1/2} \quad (\epsilon \to \epsilon_c)\]

\textbf{Proof}: The radial motion (Eq. 3.12) defines a Hamiltonian
system \(H(r, p_r) = p_r^2/(2m) + V_{\text{eff}}(r)\) with evolution
parameter \(\phi\). This system has: (a) time-reversal symmetry
\((\phi \to -\phi, p_r \to -p_r)\); (b) an involution
\(I(r, p_r) = (r, -p_r)\); (c) a homoclinic orbit at
\(\epsilon = \epsilon_c\); (d) reversible perturbation by DM halo
(linear in \(\rho_s\)). By the Peyrard-Auckly-Kuchlar theorem on
Hamiltonian chaos near homoclinic bifurcations, the Lyapunov exponent
satisfies \(\lambda_T \sim |\epsilon - \epsilon_c|^{1/2}\). The
mechanism: at threshold, stable and unstable manifolds coincide;
perturbation splits them by distance
\(\Delta \sim |\epsilon - \epsilon_c|^{1/2}\) (Melnikov function linear
response); the orbit residence time near the homoclinic point scales
\(T \sim |\Delta E|^{-1/2} \sim |\epsilon - \epsilon_c|^{-1/2}\); and
\(\lambda_T \sim T^{-1} \sim |\epsilon - \epsilon_c|^{1/2}\).
\(\square\)

\textbf{Numerical support} (Paper I Table 3, Case-III):
\(\lambda_T^2/\epsilon \approx \text{constant}\) in the chaotic regime,
consistent with \(\beta = 0.5\).

\begin{center}\rule{0.5\linewidth}{0.5pt}\end{center}

\subsection{8. Distance Equivalence
Classes}\label{distance-equivalence-classes}

\subsubsection{8.1 Unified Comparison
Table}\label{unified-comparison-table}

{\def\LTcaptype{none} % do not increment counter
\begin{longtable}[]{@{}
  >{\raggedright\arraybackslash}p{(\linewidth - 6\tabcolsep) * \real{0.2500}}
  >{\raggedright\arraybackslash}p{(\linewidth - 6\tabcolsep) * \real{0.2500}}
  >{\raggedright\arraybackslash}p{(\linewidth - 6\tabcolsep) * \real{0.2500}}
  >{\raggedright\arraybackslash}p{(\linewidth - 6\tabcolsep) * \real{0.2500}}@{}}
\toprule\noalign{}
\begin{minipage}[b]{\linewidth}\raggedright
Framework
\end{minipage} & \begin{minipage}[b]{\linewidth}\raggedright
\(\Gamma\) (source distance)
\end{minipage} & \begin{minipage}[b]{\linewidth}\raggedright
\(\delta\) (observed distance)
\end{minipage} & \begin{minipage}[b]{\linewidth}\raggedright
Critical exponent \(\alpha\)
\end{minipage} \\
\midrule\noalign{}
\endhead
\bottomrule\noalign{}
\endlastfoot
D1-grav-waves & \(d_{TV}\) & \(\zeta(\theta)\) & 1.84 (log-corrected
quadratic) \\
D16+D24 & \(d_{g_F}\) & \(D_{KL}\) & 2.00 (exact quadratic, Amari) \\
Bayesian & \(d_{g_F}\) & Fisher-Rao posterior & 2.00 \\
D26 & Pole-perturbation & \(G_{\text{amp}}\) & \(\infty\) (pole
divergence) \\
EMRI & Metric-perturbation \(\|h_{\mu\nu}\|\) & \(\lambda_T\) & 0.50
(Peyrard-Auckly) \\
DistCat & \(n/a\) (framework-level) & \(\delta = \Phi(\Gamma)\) & Varies
per example \\
\end{longtable}
}

\subsubsection{8.2 D1 vs D16+D24}\label{d1-vs-d16d24}

Both exhibit quadratic asymptotics near critical points.
\(\zeta(\theta) - 0.5 \sim \theta^2 \ln \theta\) and
\(D_{KL} \sim \frac{1}{2} g_F \epsilon^2\). They are asymptotically
equivalent:
\(D(\mathcal{B}_{\text{D1}}) \cong_{\text{asymp}} D(\mathcal{B}_{\text{D16+24}})\).

\subsubsection{8.3 D26 vs EMRI}\label{d26-vs-emri}

Both share: power-law divergence of amplification near threshold
(\(G_{\text{amp}} \sim \epsilon^{-1}\) and
\(\lambda_T \sim \epsilon^{1/2}\)); resonance-peak structure; stability
bounds (\(H_\infty\) robustness and MSS chaos bound). They belong to the
same distance equivalence class.

\subsubsection{8.4 D1 vs EMRI: GW Observational
Bridge}\label{d1-vs-emri-gw-observational-bridge}

The connection is through \textbf{spectral structure}: - Chaotic EMRI
GWs: broadband continuous spectrum; non-chaotic: discrete line spectrum
- D1 isotropic GWB: continuous HD correlation; D1 discrete SMBHB
sources: anisotropic HD correlation

Both dichotomies (continuous vs discrete) are manifestations of the same
DistCat property: the continuity/discreteness of \(\delta\) is a
function of \(\Gamma\). When \(\Gamma\) changes continuously, \(\delta\)
is continuous (non-chaotic orbits \(\to\) HD isotropic background); when
\(\Gamma\) undergoes bifurcation, \(\delta\)'s discreteness is broken
(chaotic orbits \(\to\) HD anisotropy).

This establishes
\(D(\mathcal{B}_{\text{D1}}) \cong_{\text{struct}} D(\mathcal{B}_{\text{EMRI}})\):
they share the same distance-export asymptotic pattern
\(\delta \sim \Gamma^\alpha\) (\(\alpha \in (1, 2)\)) despite different
critical exponents.

\begin{center}\rule{0.5\linewidth}{0.5pt}\end{center}

\subsection{9. Discussion}\label{discussion}

\subsubsection{\texorpdfstring{9.1 Universal Nature of
\(\delta = \Phi(\Gamma)\)}{9.1 Universal Nature of \textbackslash delta = \textbackslash Phi(\textbackslash Gamma)}}\label{universal-nature-of-delta-phigamma}

The five examples demonstrate that distance export is not an isolated
phenomenon but a \textbf{universal structural pattern} appearing in
AQFT, information geometry, Bayesian statistics, control theory, and
general relativity. The common feature is a natural transformation
\(\eta: \delta \Rightarrow \Gamma\) connecting an observable distance to
a more fundamental distance.

\subsubsection{9.2 Critical Exponents as Categorical
Invariants}\label{critical-exponents-as-categorical-invariants}

The critical exponent \(\alpha\) in \(\delta \sim \Gamma^\alpha\) is a
\textbf{2-categorical invariant} determined by the naturality of
\(\eta\) near critical points. Different \(\alpha\) values do not
contradict equivalence --- they reflect different members of the same
distance equivalence class \([\mathcal{B}]\).

\subsubsection{9.3 Implications for Theory
Unification}\label{implications-for-theory-unification}

If distinct physical theories share the same distance-export structure,
they may be unified as different 0-objects within the same 2-category.
Theorem 2.6 (observation infeasibility) implies that observational data
alone cannot distinguish between theories in the same distance
equivalence class. Theory choice requires additional criteria:
simplicity of the base category \(\mathbf{C}\), naturalness of the
structure functor \(\mathcal{F}\), or empirical adequacy beyond distance
measurements.

\subsubsection{9.4 Limitation: Discreteness of the GR
Category}\label{limitation-discreteness-of-the-gr-category}

The EMRI example (Example 5) has a \textbf{discrete source category}
\(\mathbf{C}_{\text{EMRI}}\): the Dehnen metrics are analytic in the
parameters \((\rho_s, r_s)\) and \(r\), so an isometric embedding
between two spacetimes requires parameter equality. Consequently,
1-morphisms are only identity morphisms, and all categorical claims
(naturality of \(\eta\), functoriality of morphisms, 2-morphism
composition) are \textbf{vacuously satisfied} for this specific
0-object.

This is not a flaw in the framework --- it is a \textbf{feature of the
distance export principle}. The principle \(\delta = \Phi(\Gamma)\)
holds for EMRI even when the categorical machinery is empty,
demonstrating that the principle is a property of the underlying
geometry of the parameter space, not merely a consequence of categorical
composition. The non-trivial categorical structure (non-discrete
\(\mathbf{C}\), non-vacuous naturality) is exhibited by the other four
examples: D1-grav-waves (AQFT with rich morphism structure between local
algebras), D16+D24 (statistical models with measure-preserving maps),
Bayesian inference (statistical models with stochastic kernels), and D26
(LTI systems with feedback morphisms).

Future work will explore whether \(\mathbf{C}_{\text{EMRI}}\) can be
enlarged to include \textbf{approximately isometric embeddings} (where
\(||\psi^* g_2 - g_1|| < \epsilon\)) or \textbf{local patch embeddings}
(where isometry is required only on a compact region \(U \subset M\)),
which would make the GR example genuinely non-discrete.

\subsubsection{9.5 Open Problems}\label{open-problems}

\begin{enumerate}
\def\labelenumi{\arabic{enumi}.}
\tightlist
\item
  \textbf{Rigorous 2-category proof for EMRI}: \(\square\) Completed
  (iteration 71). See \texttt{DISTCAT-EMRI-2category-proof.md}. All 12
  axioms verified: category structure, functoriality of \(\mathcal{F}\)
  and \(\mathcal{O}\), Lipschitz continuity of \(\eta\), distance export
  functor well-definedness, interchange law, and structural equivalence
  to D26.
\item
  \textbf{D1 vs EMRI quantitative equivalence}: \(\square\) Completed
  (iteration 72). D1 (α≈1.84, log-corrected quadratic) and EMRI (α=0.5,
  homoclinic bifurcation) belong to \textbf{different universality
  classes} with incommensurable exponents. D1's critical point is a
  geometric boundary (θ→0 is minimum angular separation, not a
  bifurcation), while EMRI's is a dynamical phase transition
  (KAM-breakup). They share the DistCat 2-categorical structure but not
  numerical exponent classes. See
  memory/72-d1-emri-quantitative-comparison.md.
\item
  \textbf{Extension to quantum gravity}: Can the DistCat framework
  accommodate distances derived from quantum gravitational structures
  (e.g., AdS/CFT entanglement entropy)?
\item
  \textbf{Optimal distance export}: \(\square\) Completed (iteration
  73). \texttt{DISTCAT-optimal-distance-export.md} --- 证明 Fisher-Rao
  距离是唯一最优预距离 (定理 1.4 + 定理 3.1)。所有 5 个特例都是
  Fisher-Rao 最优性的特例。关键结果：Fisher-Rao
  是唯一使数据处理的信息不等式取等号的度量 (Amari/Cencov/Chentsov
  唯一性定理)。
\end{enumerate}

\begin{center}\rule{0.5\linewidth}{0.5pt}\end{center}

\subsection{References}\label{references}

\begin{enumerate}
\def\labelenumi{\arabic{enumi}.}
\tightlist
\item
  Hellings, R. \& Downs, G. (1983). Upper bounds on the isotropic
  gravitational radiation background from pulsar timing. \emph{ApJ},
  265, L39.
\item
  Amari, S. (1985). Differential-geometrical methods in statistics.
  \emph{Springer}.
\item
  Hashimoto, K. \& Tanahashi, N. (2017). Universality in chaos of
  particle motion near black hole horizon. \emph{Phys. Rev.~D}, 95,
  024007.
\item
  Peyrard, M. \& Auckly, D. Variational characterization of the Lyapunov
  exponent near Hamiltonian chaos threshold. \emph{Physica D} (2000).
\item
  Das, S., Dalui, S., Lee, B.-H. \& Cai, Y.-F. (2025). EMRI Chaos I:
  Homoclinic orbit and near-horizon chaos. \emph{arXiv:2511.03657}.
\item
  Das, S., Dalui, S., Lee, B.-H. \& Cai, Y.-F. (2025). EMRI Chaos II:
  Chaotic imprints in gravitational waves. \emph{arXiv:2512.04848}.
\item
  Al-Badawi, A., Shaymatov, S. \& Sekhmani, Y. (2025). Schwarzschild
  black hole in galaxies surrounded by a dark matter halo. \emph{JCAP},
  02, 014.
\item
  Kehagias, A. \& Riotto, A. (2024). On the symmetry interpretation of
  the Hellings-Downs curve.
\item
  Maldacena, J., Shenker, S. \& Stanford, D. (2016). A bound on chaos.
  \emph{JHEP}, 08, 106.
\item
  Bombelli, L. \& Calzetta, E. (1992). Chaos around a black hole.
  \emph{Class. Quant. Grav.}, 9, 2573.
\item
  Wiggins, S. (2003). Introduction to Applied Nonlinear Dynamical
  Systems and Chaos. \emph{Springer}.
\item
  Gelfand, I. \& Naimark, M. (1943). On the imbedding of normed rings.
  \emph{Uspekhi Mat. Nauk}, 8, 3-33.
\end{enumerate}

\end{document}
